% --- Codificación e Idioma ---
\usepackage[utf8]{inputenc}
\usepackage[spanish, es-tabla]{babel} % Idioma español y uso de "Tabla" en lugar de "Cuadro"

% --- Matemáticas y Física ---
\usepackage{amsmath, amssymb, amsfonts} % Símbolos y entornos matemáticos
\usepackage{siunitx} % Excelente para escribir unidades con formato profesional (ej. \SI{220}{\ohm})

% --- Gráficos y Tablas ---
\usepackage{circuitikz} % Para dibujar circuitos eléctricos
\usepackage{subcaption} % Para subfiguras dentro de una figura
\usepackage{graphicx} % Para insertar imágenes
\usepackage{float} % Para forzar la posición de imágenes y tablas con [H]
\usepackage{booktabs} % Para tablas con líneas profesionales (\toprule, \midrule, \bottomrule)
\usepackage{csvsimple} % Para importar datos directamente desde archivos CSV
\usepackage{pgfplots} % Para generar gráficos a partir de datos
\usepackage{multicol} % Para crear columnas en tablas o texto
\pgfplotsset{compat=1.18} % Versión de compatibilidad para pgfplots

% --- Formato y Diseño ---
\usepackage{geometry} % Configuración de márgenes
\geometry{
    a4paper,
    top=2.5cm,
    bottom=2.5cm,
    left=2.5cm,
    right=2.5cm
}
\usepackage{fancyhdr} % Para personalizar encabezados y pies de página
\usepackage{hyperref} % Para que el índice y las referencias sean enlaces clickeables
\hypersetup{
    colorlinks=true,
    linkcolor=black,
    filecolor=magenta,
    urlcolor=blue,
    citecolor=black
}
\usepackage[bottom, hang, flushmargin]{footmisc} % Para mejorar la apariencia de las notas al pie

